\ifx\ISMAIN\undefined
    \documentclass{AsignaturaResumen}
    \usepackage{etoc}
    \SetTitulo{Diseño de experimentos}
    \setabbreviationstyle[acronym]{long-short}
\makeglossaries

    % Glosario

\newglossaryentry{latex}{
    name={latex},
    description={Procesador de textos}
}

    % Acrónimos

\newacronym{adn}{ADN}{Ácido desoxirribonucléico}
\newacronym{si}{SI}{Sistema Internacional}
\newacronym{ros}{ROS}{especies reactivas de oxígeno}
\newacronym{iubmb}{IUBMB}{Unión Internacional de Bioquímica y Biología Molecular}
    \begin{document}
    \dominitoc
    \faketableofcontents
    \chapter{Introducción a análisis estadísticos}
\fi

\section{Normalidad - Lilliefors}

\begin{align*}
    H_{0} &: \text{Los datos son normales}
    \\
    H_{a} &: \text{Los datos NO son normales}
\end{align*}

\begin{gather*}
    \bar{x} = \frac{\sum{x_{i}}}{n}
    \\\\
    \sigma = PENDIENTE
    \\\\
    f(x)_{i} = \frac{1}{n_{i}}
    \\\\
    s(x) = \sum{f(x)}
    \\\\
    z = \frac{x_{i} - \bar{x}}{s}
    \\\\
    F(x) \ \text{se obtiene de las tablas de z. Si} \ z > 1 \longrightarrow F(x) = 1 - F(x)_{\text{Tabla}}
    \\\\
    \lvert d_{i} \lvert = \lvert F(x) - S(x) \lvert
\end{gather*}


\section{Pruebas de homogeneidad}

\subsection{Bartlett}
\begin{align*}
    H_{0} &: \text{Los datos son homogéneos}
    \\
    H_{a} &: \text{Los datos NO son homogéneos}
\end{align*}

\begin{itemize}
    \item $n_{i}$ = Número de observaciones
    \item $\bar{x}$ = Promedio
    \item $V_{i}$ = Grados de libertad
\end{itemize}

\begin{gather*}
    V_{i} = n - 1
    \\\\
    \bar{x} = \frac{\sum{x_{i}}}{n}
    \\\\
    ss_{i} = \sum (x- \bar{x})^2
    \\\\
    s^{2} = \frac{ss_{i}}{V_{i}}
    \\\\
    s^{2}_{p} = \frac{\sum{ss_{i}}}{\sum{V_{i}}}
    \\\\
    B = 2.30259 \cdot  \left[ \left(\log(s_{p}^{2})\right) \cdot \left(\sum{V_{i}}\right) - \left(\sum{V_{i} \cdot \log(s^{2})}\right) \right]
    \\\\
    C = 1 + \frac{1}{3(K-1)}
    \\\\
    \chi^{2} = \frac{B}{C}
\end{gather*}

\subsection{Fisher}

Esta prueba se utiliza para analizar la homogeneidad de 2 columnas de mínimo 4 datos cada una.

En esta prueba las hipótesis son:

\begin{align*}
    H_{0} &: \text{Los datos son homogéneos}
    \\
    H_{a} &: \text{Los datos NO son homogéneos}
\end{align*}


\section{Pruebas t}

\subsubsection{Tipos de pruebas t y usos}

\begin{table}[H]
\centering
\caption{}
\label{tab:}
{\small\setstretch{1.35}\rowcolors{3}{White}{Grey}

\begin{tabular}{m{0.2\linewidth}m{0.2\linewidth}m{0.2\linewidth}m{0.2\linewidth}}\hline
    
    \textbf{Prueba t} & \textbf{Normalidad} & \textbf{Homogeneidad} & \textbf{Observaciones} \\ \hline
    
    Pareada                                      & $\checkmark$ & $\checkmark$ &    Relacionadas    \\ 
    Varianzas iguales (mismas observaciones)     & $\checkmark$ & $\checkmark$ &   Igual cantidad   \\ 
    Varianzas iguales (diferentes observaciones) & $\checkmark$ & $\checkmark$ & Diferente cantidad \\ 
    Welch                                        & $\checkmark$ & $\mathrm{X}$ &        ---         \\
    Mann-Whitney                                 & $\mathrm{X}$ &     ---      &        ---         \\ 

    \hline
\end{tabular}}
\end{table}

\subsection{Pareada}

Para datos normales, homogéneos y relacionados

\begin{equation*}
    t = \frac{\bar{d}}{s_{d} \div \sqrt{n}}
\end{equation*}


\subsection{Varianzas iguales}

\subsubsection{Observaciones iguales}

Para datos normales, homogéneos, con mismo número de observaciones ($n_{i}$)

\begin{equation*}
    t = \frac{\bar{x}_{1} - \bar{x}_{2}}{s_{p}\sqrt{\frac{1}{n_{1}} + \frac{1}{n_{2}}}}
\end{equation*}

Donde:

\begin{equation*}
    s_{p} = \sqrt{\frac{s_{1}^{2} (n_{1}-1) + s_{2}^{2} (n_{2}-1)}{n_{1} + n_{2} - 2}}
\end{equation*}


\subsubsection{Observaciones diferentes}

Para datos normales, homogéneos, con diferente número de observaciones ($n_{i}$)

\begin{equation*}
    t = \frac{\bar{x}_{1} - \bar{x}_{2}}{\sqrt{s^{2} \cdot (\frac{1}{n_{1}} + \frac{1}{n_{2}})}}
\end{equation*}

Donde:

\begin{equation*}
    s^{2} = \frac{\sum{(x_{1} - \bar{x}_{1})^{2}} + \sum{(x_{2} - \bar{x}_{2})^{2}}}{n_{1} + n_{2} - 2}
\end{equation*}


\subsection{Welch}

Para datos normales pero no homogéneos

\begin{equation*}
    t = \frac{\bar{x}_{1} - \bar{x}_{2}}{\sqrt{\frac{s_{1}^{2}}{n_{1}} + \frac{s_{2}^{2}}{n_{2}}}}
\end{equation*}

Donde:

\begin{equation*}
    \mathrm{df} = \frac{ (\frac{s_{1}^{2}}{n_{1}} + \frac{s_{2}^{2}}{n_{2}})^{2} }{ \frac{\left( \frac{s_{1}^{2}}{n_{1}} \right)^{2}}{n_{1}-1}  + \frac{\left(\frac{s_{2}^{2}}{n_{2}} \right)^{2}}{n_{2}-1} }
\end{equation*}


\subsection{Mann-Whitney}

Para datos no normales

\subsubsection{Menos de 20 datos}

\begin{gather*}
    U = n_{1} \cdot n_{2} + \frac{n_{1} \cdot (n_{1} + 1)}{2} - R_{1}
    \\\\
    U^{\prime} = n_{1} \cdot n_{2} + \frac{n_{2} \cdot (n_{2} + 1)}{2} - R_{2}
\end{gather*}

Donde $R_{n}$ se obtiene al ordenar los datos de mayor a menor y asignarles un número

\subsubsection{Más de 20 datos}

\begin{gather*}
    Z_{u} = \frac{U - \bar{U}}{U_\sigma}
    \\\\
    \bar{U} = \frac{n_{1} \cdot n_{2}}{2}
    \\\\
    U_{\sigma} = \sqrt{\frac{n_{1} \cdot n_{2} \cdot (n_{1} + n_{2} + 1)}{12}}
\end{gather*}



\section{Pruebas ANOVA}

\subsection{ANOVA paramétrico}



\subsection{ANOVA no paramétrico}



\ifx\ISMAIN\undefined
    \end{document}
\fi